\documentclass[instructions.tex]{subfiles}
\begin{document}
 
\chapter{Il faut prier pour le roi, le respecter, lui obéir ; être soumis aux magistrats.}
 
\primo Saint Paul ordonne de prier pour les princes et les rois. Nous devons prier particulierement pour celui que Dieu nous a donné, demander les bénédictions du Tout-Puissant sur le royaume, sur la famille royale et sur la personne du roi. Le plus grand bonheur d'un État, c'est d'être dans la foi catholique : prions le Seigneur qu'il conserve dans notre monarque un zèle ardent pour protéger l'Église romaine et la foi.
 
Craignez Dieu, honorez le roi, dit saint Pierre\footnote{Deum timete, regem honorificate. 1. Petr. 2. 17.}. C'est manquer de crainte de Dieu que de manquer de respect pour le roi et pour ses ordres, parce que l'autorité du prince est une émanation de l'autorité de Dieu. Ce respect doit même être intérieur. Gardez-vous, dit le Sage, de juger mal du roi, même dans votre pensée\footnote{In cogitatione tua regi ne detrahas. Eccles. 10. 20.}.
 
Que toute âme soit donc soumise aux puissances, dit saint Paul : que cette soumission soit sainte et volontaire, parce que Dieu l'ordonne, et parce qu'il est juste. Chérissons notre prince, soumettons-nous à lui, le regardant comme notre père, comme le défenseur de son peuple.
 
Refuser de payer le tribut, c'est manquer de respect et d'obéissance au roi. Jésus-Christ n'a-t-il pas payé lui-même le tribut au prince, et n'ordonne-t-il pas de rendre à César ce qui appartient à César ? N'est-ce pas le prince qui conserve nos biens, qui défend notre vie, qui protége la religion, qui empêche les assassinats et les brigandages? Quel désastre, si nous voyions des ennemis furieux ravager nos campagnes, mettre tout à feu et à sang ; saccager nos maisons et nos villes! si nous voyions des armées d'infidèles et d'hérétiques au milieu de nous, massacrer les ministres de Jésus-Christ, brûler ses églises, détruire ses autels! Quelle obligation n'avons-nous pas au roi d'empêcher tant de malheurs! N'est-il pas juste que nous l'aidions par nos facultés à soutenir les dépenses qu'il est obligé de faire pour notre sûreté et pour notre garde?
 
\secundo On doit encore obéir aux puissances, surtout dans ce qui regarde la police et le bon ordre. Les princes, les magistrats, ont fait de sages réglemens pour empêcher les abus. Tels sont les édits contre les filles et les veuves qui cèlent leur grossesse, et qui font périr leur fruit ; telles sont les ordonnances contre la fréquentation des cabarets dans le lieu du domicile, contre les profanations des saints jours, contre les scandales et les courses nocturnes, contre les désordres qui arrivent à la célébration des noces, etc.
 
Si les peuples ne s'y soumettent pas, ils sont doublement coupables, parce qu'ils résistent à une double autorité, à celle du prince et des magistrats, et à celle de l'Église. Or, ceux qui résistent à l'autorité, résistent à l'ordre de Dieu, dit saint Paul, et s'acquièrent la damnation\footnote{Qui resistit potestati, Dei ordinationi resistit ; qui autem resistunt, ipsi sibi damnationem acquirunt. Rom. 13. 2.}.
 
\section{Résolutions}
 
\primo Je respecterai le roi ; je lui demeurerai fidèle, et je prierai pour lui, pour sa famille, et pour les magistrats qu'il emploie dans le gouvernement.
 
\secundo J'obéirai à l'autorité temporelle dans tout ce qui est de l'ordre temporel, observant soigneusement cette maxime de l'Évangile : Rendez à César ce qui est à César, et à Dieu ce qui est à Dieu\footnote{Matth. 22. 21.}. 

\prayer{Mon Dieu, donnez à vos peuples fidèles des rois qui aiment la religion, protégeant l'Église, et gouvernent leurs sujets avec sagesse.}
 
\end{document}
